\documentclass[10pt,a4paper]{article}
\usepackage{fontspec,amsmath,amssymb,graphicx}
\usepackage[margin=18mm]{geometry}
\setmainfont{Times New Roman}
\setlength{\parindent}{0pt}
\setlength{\parskip}{5pt}
\allowdisplaybreaks
\sloppy
\emergencystretch=2em
\begin{document}
\section*{BT\_4\_5\_3\_04 — Lời giải bổ sung}
\subsection*{Kiến thức và trình tự giải}
\textbf{Bước 1 — Loại đạo hàm cấp một bằng tính đạo hàm đầy đủ.} Với $u_t=Ku_{xx}+Au_x+Bu$, đặt $u=e^{-hx}v$, $h=A/(2K)$. Khi đó


\[
u_t=e^{-hx}v_t,\quad u_x=e^{-hx}(v_x-hv),\quad
u_{xx}=e^{-hx}(v_{xx}-2hv_x+h^2v).
\]


Thay vào PDE, hệ số $v_x$ là $A-2Kh=0$, hệ số $v$ là $Kh^2-Ah+B=B-A^2/(4K)$. Vì thế $v_t=Kv_{xx}+Cv$, $C=B-A^2/(4K)$; có thể tiếp tục đặt $v=e^{Ct}z$ để $z_t=Kz_{xx}$.

\textbf{Bước 2 — Đổi dữ liệu và biên cùng phép đổi ẩn.} $v(x,0)=e^{hx}f(x)$. Dirichlet thuần nhất giữ nguyên; Neumann biến thành $v_x-hv=0$, nên không giữ nguyên một cơ sở cosine nếu $A\ne0$. Với Dirichlet:


\[
u=e^{-hx}\sum_{n\ge1}\left(\frac2L\int_0^Le^{hy}f(y)\sin\frac{n\pi y}Ldy\right)
e^{[C-K(n\pi/L)^2]t}\sin\frac{n\pi x}L.
\]


\textbf{Bước 3 — Tính hệ số mũ–đa thức.} Dùng nguyên hàm


\[
\int e^{hy}\sin ky\,dy
=\frac{e^{hy}(h\sin ky-k\cos ky)}{h^2+k^2}.
\]


Với $y e^{hy}$ hoặc $y(1-y)e^{hy}$, lấy đạo hàm theo $h$ của tích phân đã tính trên cận cố định; trình bày đạo hàm để tìm hệ số thực. Cách này tránh để lại một tích phân sơ cấp chưa thực hiện.

\textbf{Bước 4 — Hệ số không gian biến thiên.} Với $x^2u_{xx}$ trên $x>0$, đặt $z=\ln x$: nếu $U(z,t)=u(e^z,t)$ thì $u_x=U_z/x$, $x^2u_{xx}=U_{zz}-U_z$. Đổi miền và biên sang biến $z$ rồi dùng bước 1. Biên không thuần nhất nên tìm trạng thái dừng trước khi khai triển.

\textbf{Bước 5 — Duy nhất và ổn định.} Với hiệu nghiệm Dirichlet, $I'=-K\int w_x^2+2BI\le2BI$ suy $w=0$ khi $I(0)=0$. Giới hạn phụ thuộc các số $C-K\lambda_n$ có hệ số khác không: âm thì giảm, bằng không thì tồn tại mode dừng, dương thì tăng. Phải kiểm tra tính thống nhất giữa miền PDE và vị trí biên trước khi dùng công thức.
\subsection*{Lời giải áp dụng cho bài}
\textbf{Bước 1. Kiểm tra miền.}
Đề ghi PDE trên $(0,2)$, hai biên tại $0,1$ và dữ liệu đầu chỉ trên $(0,1)$. Vì $1$ là điểm trong của miền PDE, đây không phải bài Dirichlet chuẩn trên đoạn $(0,2)$. Cần làm rõ miền dự kiến trước khi đưa ra đáp án theo tách biến trên đoạn.
\textbf{Bước 2. Giải bài nếu miền dự kiến là $(0,1)$.}
Đặt $h=\beta/2$, $u=e^{-hx}v$. Tính $u_x=e^{-hx}(v_x-hv)$, $u_{xx}=e^{-hx}(v_{xx}-2hv_x+h^2v)$; suy $v_t=v_{xx}-h^2v$ và $v(x,0)=e^{hx}x(1-x)$.
\textbf{Bước 3. Thực hiện tích phân hệ số.}
Với $k=n\pi$, $D=h^2+k^2,E=(-1)^ne^h$, nguyên hàm cho $J=\int_0^1e^{hy}\sin ky=k(1-E)/D$. Hệ số $2\int e^{hy}y(1-y)\sin ky=2(J'-J'')$:
\[c_n=\frac{4k\{(1-h)-(1+h)E\}}{D^2}
-\frac{16kh^2(1-E)}{D^3}.\]
Đạo hàm được tính từ $E'=E,D'=2h$: $J'=-kE/D-2kh(1-E)/D^2$ và $J''=-kE/D+(4khE-2k(1-E))/D^2+8kh^2(1-E)/D^3$.
\textbf{Bước 4. Ghép và kiểm tra.}
\[u=e^{-hx}\sum_{n\ge1}c_ne^{-(n^2\pi^2+h^2)t}\sin n\pi x,\quad0<x<1.\]
Mỗi mode thỏa biên, ODE và dữ liệu biến đổi; cộng lại trả về PDE gốc. Mọi mũ âm, nên giới hạn đều $0$ với mọi $\beta$ thực. Công thức trên giải miền sửa $(0,1)$. Nếu chủ ý giữ miền PDE $(0,2)$ và yêu cầu nghiệm cổ điển qua $1$ với $t>0$, tính giải tích theo không gian buộc thác triển khác: với $v=e^{hx}u$, phản xạ lẻ qua $1$ cho $u(x,t)=-e^{2h(1-x)}u(2-x,t)$ trên $(1,2)$. Nó xác định dữ liệu và biên phía phải không được in; không được tự coi sine $n\pi x$ là cơ sở Dirichlet tại $2$.\end{document}